\documentclass[UTF8,a4paper,12pt]{ctexart}

\usepackage[a4paper,left=2.6cm,right=2.6cm,top=2.5cm,bottom=2.5cm]{geometry}
\usepackage{graphicx}
\usepackage{booktabs}
\usepackage{longtable}
\usepackage{array}
\usepackage{amsmath}
\usepackage{xcolor}
\usepackage{listings}
\usepackage{hyperref}
\usepackage{enumitem}
\usepackage{float}

\hypersetup{
  colorlinks=true,
  linkcolor=black,
  urlcolor=blue,
  citecolor=black
}

\setlength{\parindent}{2em}
\setlength{\parskip}{0.35em}
\renewcommand{\arraystretch}{1.35}
\setlist{itemsep=0.25em,topsep=0.3em}

\lstdefinelanguage{Verilog}{
  morekeywords={module,endmodule,input,output,inout,wire,reg,assign,always,begin,end,
    if,else,case,endcase,localparam,parameter,function,endfunction,task,endtask,
    posedge,negedge,or,integer,for,while,default},
  sensitive=true,
  morecomment=[l]{//},
  morecomment=[s]{/*}{*/},
  morestring=[b]"
}
\lstset{
  language=Verilog,
  basicstyle=\small\ttfamily,
  keywordstyle=\color{blue},
  commentstyle=\color{gray},
  numbers=left,
  numberstyle=\scriptsize\color{gray},
  frame=single,
  breaklines=true,
  columns=fullflexible,
  showstringspaces=false
}

\title{自动售货机控制系统设计}
\author{学生姓名：\underline{\hspace{4cm}}\quad 学号：\underline{\hspace{4cm}}\\[0.6em]
任课教师：李腾}
\date{2026 年 \hspace{1cm} 月 \hspace{1cm} 日}

\begin{document}

\begin{center}
  {\zihao{-2}\bfseries 东南大学信息科学与工程学院\\[0.8em]}
  {\zihao{2}\bfseries 数字系统课程设计实验报告}
\end{center}
\vspace{1em}

\begin{center}
\begin{tabular}{>{\raggedright\arraybackslash}p{0.25\linewidth}p{0.67\linewidth}}
题目 & 自动售货机控制系统 \\
学期 & 2026--2027（暑期学校） \\
开发平台 & Vivado 2026.1，Verilog HDL \\
目标器件 & AMD/Xilinx Artix-7 XC7A75T-2FGG484，HX7A75C 开发板 \\
实验类型 & FPGA 数字系统综合设计
\end{tabular}
\end{center}

\tableofcontents
\newpage

\section{系统设计}

\subsection{题目分析}

本设计实现一个基于 FPGA 的自动售货机控制系统。系统通过开发板上的四位拨码开关和四个按键完成商品选择、数量选择、订单确认、付款、出货以及找零；通过八位数码管显示当前状态、商品编号、数量、订单金额和剩余金额；通过 LED、蜂鸣器和出货控制脉冲给出状态提示。

在基础售货机流程的基础上，本设计加入了库存管理和管理员维护功能。每种商品具有独立的库存计数器，复位后初始库存为 5 件；成功出货时对应商品库存减 1，库存为零时禁止继续加入订单并给出缺货提示。管理员模式可以查看指定商品库存、补货和修改价格。找零采用人工操作模型：KEY3 每次退 1 元，KEY2 根据拨码选择退 5 元、10 元或 20 元，且不会超过当前剩余金额。

\subsection{总体方案}

系统由按键消抖与脉冲产生模块、售货机核心状态机、数码管动态扫描模块、顶层提示与蜂鸣器模块以及约束文件组成。顶层结构如图~\ref{fig:top} 所示。

\begin{figure}[H]
  \centering
  \fbox{\begin{minipage}{0.86\linewidth}
  \centering
  时钟与复位 $\quad\longrightarrow\quad$ 按键消抖/单脉冲 $\quad\longrightarrow\quad$ 售货机核心 FSM\\[0.7em]
  \hspace{2.7em}$\downarrow$\hspace{8em}$\downarrow$\hspace{9em}$\downarrow$\\[-0.2em]
  \hspace{1.2em}按键提示与蜂鸣器\hspace{2.0em} LED 状态指示\hspace{2.0em} 数码管动态扫描与显示
  \end{minipage}}
  \caption{自动售货机系统顶层结构}
  \label{fig:top}
\end{figure}

\subsection{状态机设计}

核心状态机采用十个状态编码，状态转换由按键脉冲、拨码值、金额和库存共同决定。

\begin{longtable}{>{\centering\arraybackslash}p{0.16\linewidth}p{0.25\linewidth}p{0.49\linewidth}}
\caption{售货机状态定义}\label{tab:states}\\
\toprule
状态编码 & 状态名称 & 状态功能 \\
\midrule
\endfirsthead
\toprule
状态编码 & 状态名称 & 状态功能 \\
\midrule
\endhead
\bottomrule
\endfoot
0 & ST\_IDLE & 待机，显示 00000000，等待开始选商品。 \\
1 & ST\_SELECT\_PRODUCT & 商品选择，拨码表示商品编号。 \\
2 & ST\_SELECT\_QTY & 数量选择，低两位拨码表示数量 1、2 或 3。 \\
3 & ST\_ORDER\_READY & 订单确认，可继续选择第二种商品或进入付款。 \\
4 & ST\_PAY & 付款，KEY1 投入 1 元，KEY2 根据拨码投入纸币，KEY3 确认付款。 \\
5 & ST\_VEND & 出货，产生出货脉冲并扣除对应库存。 \\
6 & ST\_CHANGE & 找零或退款，KEY3 退 1 元，KEY2 退 5、10 或 20 元。 \\
7 & ST\_ADMIN\_SELECT & 管理员商品选择。 \\
8 & ST\_ADMIN\_VIEW & 管理员库存查看和补货。 \\
9 & ST\_ADMIN\_PRICE & 管理员价格设置。 \\
\end{longtable}

正常购买流程为“待机 $\rightarrow$ 选择商品 $\rightarrow$ 选择数量 $\rightarrow$ 订单确认 $\rightarrow$ 付款 $\rightarrow$ 出货 $\rightarrow$ 找零/待机”。选择商品和数量时按 KEY4 可以清空未完成订单；拨码为 1111 时按 KEY4 直接复位整个系统。复位同时将所有商品库存恢复为 5 件、清除订单和金额。

管理员模式的入口为拨码 1110 并长按 KEY4 约 2 秒。管理员状态下可以选择商品，查看库存，按 KEY2 补充 1 件库存，按 KEY3 进入价格设置，拨码 0000--1111 对应价格 0--15 元，按 KEY1 保存价格，按 KEY4 退出管理员模式。

\subsection{输入输出规划}

\begin{table}[H]
\centering
\caption{系统输入信号规划}
\label{tab:inputs}
\begin{tabular}{>{\centering\arraybackslash}p{0.22\linewidth}p{0.2\linewidth}p{0.43\linewidth}}
\toprule
信号 & 类型 & 功能说明 \\
\midrule
clk & 输入 & 50 MHz 系统时钟。 \\
sw[3:0] & 输入 & 商品编号、数量、纸币面值、退款面值或管理员价格。 \\
key\_n[3:0] & 输入 & 四个低有效按键，经过消抖和脉冲化处理。 \\
reset & 内部 & 上电复位，或拨码 1111 时按 KEY4 产生的手动复位。 \\
\bottomrule
\end{tabular}
\end{table}

\begin{table}[H]
\centering
\caption{系统输出信号规划}
\label{tab:outputs}
\begin{tabular}{>{\centering\arraybackslash}p{0.22\linewidth}p{0.2\linewidth}p{0.43\linewidth}}
\toprule
信号 & 类型 & 功能说明 \\
\midrule
led[3:0] & 输出 & 显示待机、选择、付款、出货、找零及错误提示。 \\
seg[7:0] & 输出 & 数码管段选信号，显示数字和小数点。 \\
sel[7:0] & 输出 & 八位数码管位选信号，采用动态扫描。 \\
buzzer & 输出 & 按键、出货、错误等声音提示。 \\
vend\_pulse & 内部事件 & 出货控制脉冲。 \\
return\_coin\_pulse & 内部事件 & 单次退款或找零脉冲。 \\
sold\_out\_pulse & 内部事件 & 缺货提示脉冲。 \\
\bottomrule
\end{tabular}
\end{table}

\section{模块设计}

\subsection{按键消抖模块}

开发板按键为机械按键，按下和释放时可能产生多个快速跳变。\texttt{button\_conditioner} 模块对按键电平进行同步、稳定计数和边沿检测，只在确认按键稳定按下时输出一个时钟周期的脉冲。核心状态机只处理脉冲信号，因此一次按键只触发一次状态转移。

\subsection{售货机核心模块}

\texttt{vending\_machine\_core} 是系统的主要功能模块，内部包括状态寄存器、两条订单记录、16 个库存计数器、16 个价格寄存器、付款金额、找零金额和出货计时器。

每种商品的编号由四位拨码表示，编号范围为 0--15。库存数组为 16 个四位计数器，复位值为 5，最大值为 15。订单最多同时包含两种商品，每种商品数量为 1--3。总价计算公式为

\begin{equation}
  \mathrm{total\_due} = p_1q_1 + p_2q_2,
\end{equation}

其中 $p_i$ 为商品单价，$q_i$ 为购买数量。当前版本中，不同商品使用各自独立的库存判断。例如商品 A33 购买 3 件后，商品 A31 仍按照 A31 自身剩余库存判断，不会错误扣除 A33 的数量。

付款阶段 KEY1 表示投入 1 元硬币，KEY2 表示投入纸币。纸币面值由低两位拨码决定：00、01、10、11 分别对应 5、10、20、50 元。金额达到或超过应付金额后按 KEY3 确认付款，系统进入出货状态；金额不足时保持付款状态并给出提示。出货结束后，若剩余金额不为零则进入找零状态。

找零阶段的两类操作如下：KEY3 每次将余额减 1；KEY2 按低两位拨码选择退 5、10 或 20 元。若所选面值大于当前余额，则本次操作不执行，避免出现负余额。

\subsection{数码管显示模块}

\texttt{seven\_segment\_scan} 模块以较高频率轮流点亮八位数码管，使人眼观察到稳定的多位显示。显示内容按照状态进行复用。普通状态下，从左到右主要显示状态码、商品行列、数量、总价和辅助金额；管理员状态下显示管理员状态、商品编号、库存或价格。段选和位选均根据 HX7A75C 板上的低有效连接进行处理。

\subsection{顶层提示模块}

\texttt{vending\_machine\_top} 负责连接所有模块，并生成 LED、蜂鸣器和复位逻辑。KEY4 与拨码 1111 组合产生手动复位；拨码 1110 并长按 KEY4 时产生管理员进入脉冲。普通按键会产生短暂蜂鸣器提示，出货状态会持续提示，缺货、付款不足等异常情况通过 LED 闪烁和蜂鸣器提示。

\subsection{约束设计}

约束文件为 \texttt{hx7a75c\_vending\_machine.xdc}，目标器件为 \texttt{xc7a75tfgg484-2}。文件为时钟、拨码、按键、LED、数码管段选、数码管位选和蜂鸣器分配 FPGA 管脚，并设置 LVCMOS33 电平标准。板上实际验证时确认拨码开关向上为逻辑 0、向下为逻辑 1；数码管采用低有效显示。

\section{实验结果及结论}

\subsection{仿真方法和仿真文件}

仿真文件为 \texttt{sim/tb\_vending\_machine\_core.v}，采用脉冲任务模拟四个按键，使用较小的出货计时参数缩短仿真时间。测试覆盖以下内容：

\begin{enumerate}
  \item 复位后所有商品库存均为 5；
  \item 商品选择、数量选择和总价计算；
  \item 两种不同商品分别购买 3 件时的独立库存判断；
  \item 付款、出货、20 元和 10 元大额退款以及 1 元退款；
  \item 出货后库存递减；
  \item 管理员查看库存、补货和修改价格；
  \item 商品库存耗尽后的缺货拒绝逻辑。
\end{enumerate}

原有功能仿真测试输出为：\texttt{PASS: inventory, admin and manual change simulation completed.} 针对本次库存判断问题，测试平台已增加“A33×3 后选择 A31×3”的回归测试，要求两种商品均可加入订单且总价按各自单价计算。重新综合生成 bitstream 后，还应在开发板上完成该用例的实测确认。

\subsection{开发板验证}

实验使用 HX7A75C 开发板，通过 Vivado Hardware Manager 下载生成的 \texttt{vending\_machine\_top.bit} 文件。上电后系统自动复位；也可以将四个拨码全部拨上为 1111，再按 KEY4 进行手动复位。验证过程中观察到数码管、LED、蜂鸣器和按键提示均能正常工作。

板上主要操作与预期结果如下：

\begin{table}[H]
\centering
\caption{开发板功能验证项目}
\label{tab:boardtest}
\begin{tabular}{p{0.28\linewidth}p{0.27\linewidth}p{0.31\linewidth}}
\toprule
测试项目 & 操作 & 预期现象 \\
\midrule
系统复位 & SW=1111，按 KEY4 & 回到待机，显示 00000000，库存恢复为 5。 \\
商品选择 & 设置商品拨码，按 KEY1 & 进入商品选择或数量选择状态。 \\
数量确认 & 设置低两位拨码，按 KEY1 & 数量锁存，进入订单确认。 \\
付款 & KEY1 投 1 元或 KEY2 投纸币 & 已付款金额增加，蜂鸣器短响。 \\
确认付款 & 金额足够后按 KEY3 & 进入出货状态，LED 和蜂鸣器提示。 \\
找零 & KEY3 或 KEY2 配合拨码 & 分别退 1 元或退 5、10、20 元。 \\
管理员 & SW=1110，长按 KEY4 & 进入管理员状态，可查看、补货和改价。 \\
\bottomrule
\end{tabular}
\end{table}

\subsection{实验结论}

本设计完成了一个具有完整购买流程的 FPGA 自动售货机控制系统，实现了状态机控制、订单管理、付款结算、人工找零、动态数码管显示、LED 与蜂鸣器提示、库存管理以及管理员维护功能。通过仿真和开发板验证，系统的基本功能能够按照设计要求运行。库存判断回归测试进一步保证了不同商品之间的库存相互独立，解决了购买 A33×3 后无法继续购买 A31×3 的问题。

\section{设计特色及问题解决}

\subsection{设计特色}

本设计的主要特色包括：

\begin{enumerate}
  \item 采用明确的多状态 FSM，将选择、付款、出货、找零和管理员操作分开，便于调试和扩展；
  \item 每种商品使用独立库存计数器，成功出货才扣库存，避免选购或取消订单造成库存错误；
  \item 同时支持 16 种商品、最多两种商品组合购买和数量选择；
  \item 具备库存查看、补货和改价功能，管理员模式通过特殊拨码与长按按键进入；
  \item 找零操作同时支持 1 元和 5、10、20 元面值，能够模拟不同退款操作；
  \item 通过单脉冲按键接口、上电复位和明确的手动复位组合，提高了硬件操作稳定性。
\end{enumerate}

\subsection{开发过程中遇到的问题}

\paragraph{数码管显示异常。}
初始下载后曾出现数码管显示 88888888、部分位不亮或小数点异常的现象。分析后确认需要同时核对段选有效电平、位选有效电平、位序和约束管脚，不能只根据逻辑代码判断。通过板级诊断工程逐项测试段选和位选，最终确认开发板的数码管采用低有效控制，并修正了显示扫描的位序映射。

\paragraph{拨码和按键有效电平。}
开发板实测得到拨码向上为逻辑 0、向下为逻辑 1；DIG/LED 开关向下时数码管有效，向上时 LED 模式有效；BUZ/SW 开关向上时允许蜂鸣器工作。根据实测结果调整了操作说明，并将复位组合定义为拨码 1111 加 KEY4。

\paragraph{付款确认与退款操作混淆。}
为了避免金额刚达到商品价格时自动出货造成误操作，付款阶段将 KEY3 定义为明确的付款确认键。KEY1 和 KEY2 只负责增加付款金额；出货后进入找零状态，KEY3 逐元退款，KEY2 选择较大面额退款。

\paragraph{不同商品库存相互影响。}
原库存判断将已经选中的其他商品数量错误地加入当前商品的库存需求。例如 A33×3 后选择 A31×3 时，错误判断变成 $5\geq 3+3$，从而拒绝订单。修正后，库存判断只检查当前商品对应的库存计数器，即 $inventory[code]\geq quantity$，不同商品分别独立计算。

\paragraph{复位和状态残留。}
售货机包含库存、订单、金额和管理员状态等多个寄存器。若上一次操作停留在付款或退款状态，直接开始下一次测试可能造成误判。因此在板上测试每组流程前使用 SW=1111 加 KEY4 复位，确保状态回到待机、订单清空、库存恢复为 5。

\section{人员任务分配情况}

本项目按单人完成填写时，可将本节修改为实际情况。若为多人合作，应根据实际贡献调整百分比。

\begin{table}[H]
\centering
\caption{人员分工表}
\label{tab:members}
\begin{tabular}{p{0.24\linewidth}p{0.48\linewidth}p{0.16\linewidth}}
\toprule
成员 & 主要工作 & 贡献比例 \\
\midrule
学生姓名 & 需求分析、状态机设计、Verilog 编码、仿真、Vivado 综合实现、HX7A75C 板级调试及报告撰写 & 100\% \\
\bottomrule
\end{tabular}
\end{table}

\section{其他}

系统当前价格表由核心模块初始化。价格和库存均为易失性寄存器，手动复位后恢复到默认值；管理员改价和补货功能用于当前运行周期内的维护，不提供断电保存。若后续需要做成实际产品，可增加非易失存储器、真实硬币传感器、电机驱动模块、库存传感器和更完善的找零策略。

后续可改进内容包括：将商品编号和数量改为独立按键输入，增加订单取消后的完整退款显示，加入硬币库存约束，加入更细致的错误码显示，以及对管理员模式增加密码或专用拨码保护。

\section{附件清单}

设计源文件和约束文件建议按以下目录打包提交：

\begin{longtable}{p{0.55\linewidth}p{0.31\linewidth}}
\caption{提交附件文件列表}\label{tab:files}\\
\toprule
文件或目录 & 用途 \\
\midrule
\endfirsthead
\toprule
文件或目录 & 用途 \\
\midrule
\endhead
\bottomrule
\endfoot
\texttt{src/vending\_machine\_core.v} & 售货机核心状态机、订单、库存、付款和退款逻辑。 \\
\texttt{src/vending\_machine\_top.v} & 顶层连接、复位、LED、蜂鸣器和管理员入口。 \\
\texttt{src/button\_conditioner.v} & 按键同步、消抖和脉冲产生。 \\
\texttt{src/seven\_segment\_scan.v} & 八位数码管动态扫描与显示编码。 \\
\texttt{constraints/hx7a75c\_vending\_machine.xdc} & HX7A75C 开发板管脚和时钟约束。 \\
\texttt{sim/tb\_vending\_machine\_core.v} & 核心模块仿真测试平台。 \\
\texttt{scripts/create\_project.tcl} & 自动创建 Vivado 工程。 \\
\texttt{scripts/run\_sim.tcl} & 自动启动仿真。 \\
\texttt{scripts/run\_build.tcl} & 自动综合、实现并生成 bitstream。 \\
\texttt{report/report\_vending\_machine.tex} & 本实验报告 LaTeX 源文件。 \\
\bottomrule
\end{longtable}

\end{document}
